% ====================================================================
% Ch6(피팅 장) 통합용 추출 재료 — Chapter 2 에서 이관한 "식별성과 필요한 실험 제약"
%   원칙: 식별성·실험 프로토콜(피팅 방법론)은 물리 장(Ch1~5)이 아니라 Ch6(구 Appendix B)에 둔다(F.9b).
%   추후 Archive_old/graphite_ica_ch1_appendixB_fitting.tex 의 sec:ch6_identif 등과 병합.
%   참조 라벨(eq:ch2_*, \S\ref{sec:ch2_*})은 통합 시 재지정 필요.
%   출처: graphite_ica_ch2.tex §12 (2026-06-07 추출), 발열 챕터의 열원 식별성.
% ====================================================================
\section{발열 파라미터의 식별성과 필요한 실험 제약 (Ch2 열원 계층)}\label{sec:ch6_ch2_ident}
물리 우선 모델이라 해서 식별성(identifiability) 문제가 사라지지는 않는다. 발열 항에는 서로 다른 물리
기원이 같은 관측(전압·표면 온도)에 거의 같은 방식으로 들어가는 상관 쌍이 있어, 단일 데이터셋의 동시 자유
피팅으로는 분리되지 않는다. 해결책은 임의 clipping 이 아니라 독립 실험과 물리적 prior 로 한 겹씩 벗기는
것이다(Chapter 1 의 비순환 식별 논리와 같은 원리). 주요 상관 쌍과 그 해소 경로:

\begin{itemize}
\item $R_n$ 과 $R_{n,\eff}$: apparent 전압축 이동 계수 대 실제 dissipated heat 저항. pulse/current-interrupt/impedance
  로 ohmic·charge-transfer 성분을 분리.
\item $(\partial V_{n,\OCV}/\partial T)_q$ 와 $\Delta S_j$: 둘 다 가역열을 설명하므로 동시에 자유로 두면 double
  counting. 저율 OCV series 가 좌변 온도계수를 고정하면 우변 $\{h_\bg,\Delta S_j\}$ 합이 제약된다(Ch2 정합식).
\item $\Delta S_j$(reaction)와 $\Delta S_{a,j}$(activation): 전자는 평형 entropy($\partial U_j/\partial T$),
  후자는 rate prefactor 의 별개 양 — 대입 금지.
\item $k_{j,\mathrm{act}}$ 와 $k_{j,\lim}$: 전위에 의한 활성화 가속과 물리적 율속 병목의 분리(Chapter 1 계승; transport 진단).
\item $Q_\bg$ 와 $h_\bg$: 배경 용량 저장소와 그 entropy coefficient. 저속 OCV 의 SOC 분해 + progress/coordinate-shift 분리로 해소.
\item $h_T A$ 와 내부 열원: 표면 온도만 측정하면 heat generation 과 heat loss 가 상호 보상해 분리 불가 — 독립 열손실 교정 또는 calorimetry 필요.
\end{itemize}
필요한 독립 실험 제약(순차):
\begin{enumerate}
\item 저율 OCV/quasi-OCV 온도 series: $(\partial V_{n,\OCV}/\partial T)_q$, $U_j(T)$, $w_j(T)$, reaction entropy $\partial U_j/\partial T$ 를 평형 분기에서 제약.
\item rate series: $k_{j,\mathrm{act}}$, $k_{j,\lim}$, transport limitation, tail broadening 분리.
\item pulse/current-interrupt/impedance: ohmic 대 charge-transfer($R_n$ vs $R_{n,\eff}$).
\item 휴지 완화 자료: $\xi_j$ 완화와 thermal relaxation 분리.
\item calorimetry/온도 측정: heat source scale·가역/비가역 분리, heat loss($h_T A$) 교정, 비가역열 tail 의 forward-prediction 검정.
\end{enumerate}
한계: calorimetry 자료가 없으면 가역열 대 비가역열 분해가 전압 자료만으로는 닫히지 않으며, $\partial U_j/\partial T$ 의
staging별 분해도 측정된 anode entropy profile 의 SOC 분리를 전제로만 위계화된다.
